\subsection{Genuine geometric integrality over an algebraically closed field}
\label{subsec:geometric-domain-base-change}

\begin{theorem}[Actual rational specialization avoiding a prescribed element]
\label{algebra:rational-specialization-nonzero}
\lean{EnvelopeProperty.exists_algHom_to_algClosed,
EnvelopeProperty.exists_algHom_to_algClosed_nonzero}
\leanok
Let $k$ be an algebraically closed field. Every nontrivial finite-type
commutative $k$-algebra has an actual $k$-algebra homomorphism to $k$.
If the algebra is a domain and $a\ne0$, one can choose this actual
homomorphism so that the entire original element $a$ remains nonzero.
\end{theorem}
\begin{proof}
Choose an actual maximal ideal. Its whole quotient field is a finite
$k$-module by the Jacobson finite-type field theorem, hence algebraic.
The genuine algebraic embedding into the algebraically closed field $k$,
composed with the whole quotient map, is the required rational evaluator.
For a domain and $a\ne0$, first form the true localization at $a$.
It is a nontrivial finite-type $k$-algebra. Apply the preceding evaluator
and restrict along the actual localization homomorphism. The image of
$a$ is a unit, so its evaluated image is nonzero. No rational-point
existence certificate is a hypothesis of this construction.
\end{proof}

\begin{theorem}[Whole tensor specialization and simultaneous finite descent]
\label{algebra:whole-tensor-rational-specialization}
\lean{EnvelopeProperty.basis_baseChange_repr_tensorMap,
EnvelopeProperty.exists_tensorSpecialization_nonzero,
EnvelopeProperty.tensorMap_left_injective,
EnvelopeProperty.exists_fgSubalgebra_tensor_pair}
\leanok
\uses{algebra:rational-specialization-nonzero}
Let $k$ be a field. Every whole basis coefficient of an algebraic tensor
commutes with an actual coefficient-algebra homomorphism; an injective
coefficient homomorphism induces an injective whole tensor map.
Two elements of $F\otimes_k A$, for arbitrary commutative $k$-algebras
$F,A$, descend together to $B\otimes_k A$ for one actual finitely
generated coefficient subalgebra $B\subseteq F$.
If $k$ is algebraically closed, $B$ is a finite-type domain and both
elements of $B\otimes_k A$ are nonzero, one actual rational evaluator
$B\to k$ retains both whole tensor elements as nonzero.
\end{theorem}
\begin{proof}
Choose a genuine $k$-basis of $A$ and its true basis after coefficient
extension. Tensor induction proves that every basis coefficient of the
whole mapped element is the actual coefficient homomorphism applied
to its original coefficient. Injectivity follows coefficient by
coefficient and then by the actual basis representation inverse law.

Each original tensor element descends to an actual finitely generated
coefficient subalgebra by the tensor finite-descent theorem. The true
supremum of these two subalgebras is finitely generated. Map both original
representatives into it via their actual inclusions. Whole tensor-map
composition identifies their images with the two original elements.

For two nonzero elements, the faithful basis representation gives one
nonzero coefficient in each. Their product is nonzero in the domain $B$.
Apply the preceding rational evaluator avoiding this actual product.
Both chosen coefficients retain nonzero images. The proved whole
coefficient formula then makes both entire specialized tensors nonzero;
no support comparison is used as a substitute for pointwise equality.
\end{proof}

\begin{theorem}[Domain preservation under arbitrary domain coefficient extension]
\label{algebra:geometric-domain-base-change}
\lean{EnvelopeProperty.tensorProduct_isDomain_of_algClosed}
\leanok
\uses{algebra:whole-tensor-rational-specialization}
Let $k$ be an algebraically closed field and let $F,A$ be arbitrary
commutative domain $k$-algebras. Then the actual algebraic tensor ring
$F\otimes_k A$ is a domain. In particular, every domain $k$-algebra
remains a domain after every genuine field extension $F/k$.
There is no finite-type, completeness, separability or characteristic
restriction on either original domain.
\end{theorem}
\begin{proof}
The actual inclusion of $A$ into the tensor ring is injective over the
field $k$, so the tensor ring is nontrivial. Suppose two nonzero entire
elements have product zero. Descend both to the preceding actual finite
coefficient subalgebra $B\subseteq F$, with whole tensor representatives.
It is a domain because $F$ is. The actual inclusion tensor map is injective
by the whole basis coefficient theorem; hence the product of the two
representatives is zero in $B\otimes_k A$, while each remains nonzero.

Choose the preceding one actual rational evaluator that retains both
whole representatives. Their images still have zero product. The genuine
canonical tensor isomorphism $k\otimes_k A\simeq A$ preserves their
nonzero values and product. This contradicts the domain property of $A$.
Thus the actual tensor ring has no zero divisors, and its already proved
nontriviality makes it a domain. This is a derived algebraic producer;
geometric integrality is not an input certificate. Normality and
smoothness after coefficient extension remain separate properties.
\end{proof}
